% Transition Graphs
\section{State Transition Graphs}

\subsection{Graphical Representation}

The operational system can be visualized as a directed graph:

\subsubsection{Graph Definition}
\[
G = (V, E)
\]
where:
\begin{itemize}
    \item $V = Q$ (vertices are states)
    \item $E \subseteq Q \times \Sigma \times Q$ (labeled edges are transitions)
\end{itemize}

\subsection{Primary Transition Graph}

\subsubsection{States}
\begin{itemize}
    \item $q_0$: Initial (empty inscription)
    \item $q_1$: Instantiated (entity present)
    \item $q_2$: Counted (quantity specified)
    \item $q_3$: Related (relation established)
    \item $q_4$: Qualified (attributes added)
    \item $q_5$: Closed (validation complete) [accepting]
\end{itemize}

\subsubsection{Major Transitions}
\begin{align*}
q_0 &\xrightarrow{\text{INST}} q_1 \\
q_1 &\xrightarrow{\text{COUNT}} q_2 \\
q_1 &\xrightarrow{\text{REL}} q_3 \\
q_1 &\xrightarrow{\text{QUAL}} q_4 \\
q_1 &\xrightarrow{\text{CLOSE}} q_5 \\
q_2 &\xrightarrow{\text{CLOSE}} q_5 \\
q_3 &\xrightarrow{\text{CLOSE}} q_5 \\
q_4 &\xrightarrow{\text{CLOSE}} q_5
\end{align*}

\subsection{Reachability Analysis}

\subsubsection{Reachable States}
All states $q_1$ through $q_5$ are reachable from $q_0$.

\subsubsection{Accepting States}
States from which acceptance is possible:
\[
\text{Reachable}(F) = \{q_1, q_2, q_3, q_4, q_5\}
\]

Most states can either close immediately or continue with operations.

\subsection{Subgraph Analysis}

\subsubsection{Strongly Connected Components}
The graph is acyclic (no SCCs except individual vertices), confirming:
\begin{itemize}
    \item No recursive structures
    \item No loops or repetition of state sequences
    \item Linear progression through inscription
\end{itemize}

\subsection{Path Enumeration}

\subsubsection{Simple Paths}
Major paths from $q_0$ to accepting states:
\begin{enumerate}
    \item $q_0 \rightarrow q_1$ (entity alone)
    \item $q_0 \rightarrow q_1 \rightarrow q_2$ (entity + count)
    \item $q_0 \rightarrow q_1 \rightarrow q_5$ (entity + closure)
    \item $q_0 \rightarrow q_1 \rightarrow q_2 \rightarrow q_5$ (entity + count + closure)
\end{enumerate}

\subsubsection{Complex Paths}
Including qualifications and relations:
\begin{itemize}
    \item $q_0 \rightarrow q_1 \rightarrow q_4 \rightarrow q_5$ (qualified entity)
    \item $q_0 \rightarrow q_1 \rightarrow q_3 \rightarrow q_5$ (related entity)
    \item Multiple operations before closure
\end{itemize}

\subsection{Graph Properties}

\subsubsection{Structural Characteristics}
\begin{itemize}
    \item DAG (Directed Acyclic Graph)
    \item Single source ($q_0$)
    \item Multiple sinks (accepting states)
    \item Bounded path length (by state count)
\end{itemize}

\subsection{Visualization}

A standard graphical representation would show:
\begin{itemize}
    \item Circles for states
    \item Arrows for transitions (labeled with sign types)
    \item Double circles for accepting states
    \item Initial state marked with incoming arrow
\end{itemize}

This visualization makes the operational structure immediately apparent.
